% ============================================================================
% Fichier  : partie1/P01_C05_droit.tex
% Titre    : Obligations Juridiques et Responsabilités
% Auteur   : MINKA MI NGUIDJOÏ Thierry Emmanuel
% Version  : 1.0 -- novembre 2026
% ============================================================================
% Sources : loi n° 2010/012 du 21 décembre 2010 (texte officiel : art. 4
% point 9, 13, 26, 29, 30, 32) ; loi n° 2024/017 (principes seulement).
% ============================================================================

\chapter{Obligations Juridiques et Responsabilités}
\label{chap:droit}
\chaptermark{Obligations juridiques}

\epigraph{%
    « La loi ne demande pas d'éviter tout incident~;\\
    elle demande d'avoir fait ce qu'il fallait. »%
}{Principe du cours}

\niveauChapitre{Intermédiaire}{Chapitres~\ref{chap:gouvernance}
et~\ref{chap:cadres}~; cours SEC3266, chapitre sur le droit.}

\begin{boxobjectifs}
\begin{itemize}
    \item Citer les obligations de gestion des risques, de conservation et
          d'audit que la loi n°~2010/012 impose aux exploitants.
    \item Distinguer obligation de moyens et obligation de résultat.
    \item Identifier les responsabilités contractuelle, civile et pénale en
          jeu après un incident.
    \item Préparer les décisions juridiques d'une crise~: plainte, rançon,
          données personnelles, assurance.
    \item Constituer la preuve qu'on a fait ce qu'il fallait.
\end{itemize}
\end{boxobjectifs}

\bigskip

% ============================================================================
\section*{Question directrice : LiZbiZ avait-elle respecté la loi avant l'attaque~?}
% ============================================================================

Le 7~mars, Me~TSALA pose trois questions à M.~ETEME~: « Aviez-vous un
dispositif normalisé de gestion des risques~? Avez-vous conservé vos
journaux~? Quand a eu lieu votre dernier audit~? » Aux trois, la réponse
est embarrassée. Ce chapitre présente ce que la loi exige, et comment le
prouver.

\medskip
\textbf{Plan du chapitre.}
\begin{enumerate}
    \item Gérer les risques~: une obligation légale (section~\ref{sec:c05-risques}).
    \item Conserver et surveiller (section~\ref{sec:c05-conserver}).
    \item L'audit de sécurité (section~\ref{sec:c05-audit}).
    \item Les responsabilités en jeu (section~\ref{sec:c05-responsabilites}).
    \item Les décisions juridiques d'une crise (section~\ref{sec:c05-crise}).
\end{enumerate}

\separateur

% ============================================================================
\section{Gérer les Risques : une Obligation Légale}
\label{sec:c05-risques}
% ============================================================================

\begin{boxjuris}
\textbf{\loicam{n°~2010/012} du 21~décembre 2010 relative à la
cybersécurité et la cybercriminalité au Cameroun} (texte officiel)

\medskip
\textbf{Article~26, alinéa~1}~: \textit{« Les exploitants des systèmes
d'information prennent toutes les mesures techniques et administratives
afin de garantir la sécurité des services offerts. A cet effet, ils se
dotent de systèmes normalisés leur permettant d'identifier, d'évaluer, de
traiter et de gérer continûment les risques liés à la sécurité des
systèmes d'information dans le cadre des services offerts directement ou
indirectement. »}

\medskip
\textbf{Article~30, alinéa~1}~: \textit{« Les exploitants des systèmes
d'information évaluent, révisent leurs systèmes de sécurité et
introduisent en cas de nécessité les modifications appropriées dans leurs
pratiques, mesures et techniques de sécurité en fonction de l'évolution
des technologies. »}

\medskip
\textbf{Portée.} La loi n'impose pas une méthode, mais un \emph{système
normalisé} qui couvre les quatre temps du processus de l'ISO~31000~:
identifier, évaluer, traiter, gérer continûment. Un registre établi après
l'attaque ne prouve pas qu'on gérait les risques avant.
\end{boxjuris}

\begin{boxdef}
\termecle{Obligation de moyens}~: obligation de mettre en œuvre tous les
moyens raisonnables pour atteindre un but, sans garantir le résultat.
\termecle{Obligation de résultat}~: obligation d'atteindre le résultat
promis. La sécurité des systèmes relève, en règle générale, de la
première~: on ne reproche pas l'attaque, mais l'absence de mesures.
\end{boxdef}

% ============================================================================
\section{Conserver et Surveiller}
\label{sec:c05-conserver}
% ============================================================================

\begin{boxjuris}
\textbf{\loicam{n°~2010/012}, article~29} (texte officiel)

\medskip
\textbf{Alinéa~1}~: \textit{« Les exploitants des systèmes d'information
ont l'obligation de conserver les données de connexion et de trafic de
leurs systèmes d'information pendant une période de dix (10) ans. »}

\medskip
\textbf{Alinéa~2}~: \textit{« Les exploitants des systèmes d'information
sont tenus d'installer des mécanismes de surveillance de contrôle d'accès
aux données de leurs systèmes d'information. Les données conservées
peuvent être accessibles lors des investigations judiciaires. »}

\medskip
\textbf{Portée.} Les journaux sont à la fois une obligation légale, une
condition de la détection et la matière de l'enquête. La messagerie de
LiZbiZ ne conservait que sept jours de journaux~: l'obligation n'était
pas tenue, et l'enquête en a été privée.
\end{boxjuris}

% ============================================================================
\section{L'Audit de Sécurité}
\label{sec:c05-audit}
% ============================================================================

\begin{boxjuris}
\textbf{\loicam{n°~2010/012}} (texte officiel)

\medskip
\textbf{Article~4, point~9}~: \textit{« Audit de sécurité~: examen
méthodique des composantes et des acteurs de la sécurité, de la politique,
des mesures, des solutions, des procédures et des moyens mis en oeuvre
par une organisation, pour sécuriser son environnement, effectuer des
contrôles de conformité, des contrôles d'évaluation de l'adéquation des
moyens (organisationnels, techniques, humains, financiers) investis au
regard des risques encourus, d'optimisation, de rationalité et de
performance. »}

\medskip
\textbf{Article~13, alinéa~1}~: \textit{« Sont soumis à un audit de
sécurité obligatoire, les réseaux de communications électroniques et les
systèmes d'information, des opérateurs, des autorités de certification et
des fournisseurs de services de communications électroniques. »}

\medskip
\textbf{Article~32, alinéas~1 et~2}~: \textit{« Les réseaux de
communications électroniques et les systèmes d'information sont soumis à
un régime d'audit de sécurité obligatoire et périodique de leurs systèmes
de sécurité par l'Agence. »} \textit{« L'audit de sécurité et les mesures
d'impact de gravité sont effectués chaque année ou lorsque les
circonstances l'exigent. »}

\medskip
\textbf{Portée.} La définition légale rattache l'audit aux \emph{risques
encourus}~: auditer, c'est vérifier que les moyens sont adéquats au
regard des risques. Les conditions et les modalités de l'audit obligatoire
sont fixées par voie réglementaire (article~13, alinéa~2)~; elles se
vérifient dans les textes d'application et auprès de l'ANTIC.
\end{boxjuris}

% ============================================================================
\section{Les Responsabilités en Jeu}
\label{sec:c05-responsabilites}
% ============================================================================

\begin{tableaucours}{L{3.2cm}L{4.4cm}L{4.6cm}}{%
    \tableauHeader{Responsabilité} &
    \tableauHeader{Fondement} &
    \tableauHeader{Chez LiZbiZ}
}
Contractuelle & Inexécution d'un contrat & Pénalités versées aux clients
(74~millions)~; recours contre le prestataire \\
Civile & Dommage causé à un tiers par une faute & Clients dont les
données ont été publiées le 19~mars \\
Pénale de l'attaquant & Infractions de la loi n°~2010/012 & Plainte du
7~mars \\
Des dirigeants & Faute de gestion, devant la société et les tiers selon
le droit des sociétés & Question posée par les actionnaires \\
Administrative & Contrôle de l'autorité de régulation & Conformité à
l'article~26 \\
\end{tableaucours}
\vspace{-0.8em}
\begin{center}{\small\itshape Tableau~5.1. Les responsabilités après un incident}\end{center}

\begin{boxalert}
\textbf{Violation typique.} Croire que la qualité de victime protège de
toute responsabilité.

\smallskip
\textbf{Règle~:} l'entreprise attaquée est victime au pénal, mais peut
être responsable au civil et au contrat si elle n'avait pas pris les
mesures raisonnables. Sa meilleure défense est un dossier qui montre ce
qu'elle avait fait, avant l'incident.
\end{boxalert}

\begin{boxinfo}
\textbf{Les données personnelles.} La loi n°~2024/017 relative à la
protection des données à caractère personnel fixe le cadre camerounais
(obligations des responsables de traitement, droits des personnes,
autorité de protection). Le cours n'en cite pas les articles~: ils doivent
être lus dans le texte publié au Journal officiel, avec le conseil
juridique, avant toute décision, notamment sur l'information des personnes
concernées après une fuite.
\end{boxinfo}

% ============================================================================
\section{Les Décisions Juridiques d'une Crise}
\label{sec:c05-crise}
% ============================================================================

\begin{boxprocedure}
\textbf{Les décisions juridiques des premières 72 heures}
\begin{enumerate}
    \item \textbf{Désigner le conseil juridique} dès l'ouverture de la
          cellule de crise.
    \item \textbf{Préserver les preuves} (journaux, images, message de
          rançon) avant toute remise en état.
    \item \textbf{Déposer plainte} auprès de la police judiciaire, de la
          gendarmerie ou du procureur.
    \item \textbf{Déclarer le sinistre} à l'assureur dans le délai de la
          police.
    \item \textbf{Examiner la rançon} avec le conseil~: aucune garantie de
          restitution, risque de financer une activité criminelle, et de
          sanctions si le destinataire est visé par des mesures
          internationales.
    \item \textbf{Évaluer les obligations envers les personnes} dont les
          données sont touchées (loi n°~2024/017) et envers les clients
          (contrats).
    \item \textbf{Notifier le prestataire} et activer les clauses du
          contrat.
\end{enumerate}
\end{boxprocedure}

\begin{boxscenario}{Opération GAROUA. Le dossier de Me TSALA}
\textbf{Constatations.} Pas de registre des risques avant mars~; journaux
de la messagerie conservés sept jours~; dernier audit en 2023, hors
prestataire~; contrat sans clause de sécurité~; police d'assurance
excluant les cyberattaques.

\medskip
\textbf{Analyse} (confiance moyenne). LiZbiZ peut difficilement démontrer
qu'elle disposait du « système normalisé » de l'article~26. Son recours
contre le prestataire est affaibli par l'absence de clause. Les actions
engagées depuis avril (registre, audit, contrat révisé) ne réparent pas le
passé, mais constituent le dossier de la prochaine fois.
\end{boxscenario}

\begin{boxmetier}{Auditeur des systèmes d'information}
\textbf{Relier chaque constat à son fondement.} Dans un rapport, l'auditeur
cite pour chaque écart l'obligation qu'il méconnaît (article de loi,
clause de contrat, politique interne). Un constat sans fondement est une
opinion.
\end{boxmetier}

\separateur

\begin{boxresume}
\textbf{Ce qu'il faut retenir de ce chapitre~:}
\begin{itemize}
    \item \textbf{L'article~26} impose un système normalisé de gestion
          continue des risques~; \textbf{l'article~30}, la révision.
    \item \textbf{L'article~29} impose dix ans de conservation des données
          de connexion et une surveillance des accès.
    \item \textbf{L'audit de sécurité} est défini par la loi au regard des
          risques encourus (articles~4, 13, 32).
    \item \textbf{La victime} peut être responsable au contrat et au civil~;
          on se défend par un dossier constitué avant.
    \item \textbf{En crise}~: conseil, preuves, plainte, assureur, rançon,
          personnes concernées, prestataire.
\end{itemize}
\end{boxresume}

\begin{boxinfo}
\textbf{Points de contrôle de l'élève.} Je sais~: citer les articles~26,
29, 30 et 32 et en tirer la portée~; distinguer obligation de moyens et
de résultat~; nommer les responsabilités en jeu~; dérouler les décisions
juridiques des 72 premières heures.
\end{boxinfo}

\bigskip

\begin{boxexo}[title={\ding{45}\quad Exercice~5.1. L'article 26 \niveaubase}]
Quels sont les quatre verbes de l'article~26 qui décrivent la gestion des
risques~? À quelles étapes de l'ISO~31000 correspondent-ils~?
\end{boxexo}

\begin{boxexo}[title={\ding{45}\quad Exercice~5.2. Dix ans \niveaubase}]
Pendant combien de temps LiZbiZ doit-elle conserver les données de
connexion de sa passerelle~? Quel article le prévoit~?
\end{boxexo}

\begin{boxexo}[title={\ding{45}\quad Exercice~5.3. Moyens ou résultat \niveaubase}]
L'obligation de sécurité d'un exploitant est-elle de moyens ou de
résultat~? Quelle conséquence pour sa défense après une attaque~?
\end{boxexo}

\begin{boxexo}[title={\ding{45}\quad Exercice~5.4. La définition \niveaumoyen}]
Selon l'article~4, point~9, au regard de quoi l'audit évalue-t-il
l'adéquation des moyens~? Qu'en déduisez-vous pour le programme d'audit~?
\end{boxexo}

\begin{boxexo}[title={\ding{45}\quad Exercice~5.5. Responsabilités \niveaumoyen}]
Pour chaque responsabilité du Tableau~5.1, donnez la preuve que LiZbiZ
aurait dû pouvoir produire.
\end{boxexo}

\begin{boxexo}[title={\ding{45}\quad Exercice~5.6. La rançon \niveaumoyen}]
Rédigez la note de Me~TSALA à M.~ETEME sur la rançon~: faits, risques,
options, recommandation.
\end{boxexo}

\begin{boxexo}[title={\ding{45}\quad Exercice~5.7. Le contrat \niveaumoyen}]
Quelles clauses auraient permis à LiZbiZ de se retourner contre SAHEL
INFOGÉRANCE~?
\end{boxexo}

\begin{boxexo}[title={\ding{45}\quad Exercice~5.8. Données personnelles \niveaumoyen}]
Le 19~mars, des données de clients sont publiées. Quelles questions faut-il
poser au conseil juridique, au regard de la loi n°~2024/017~?
\end{boxexo}

\begin{boxexo}[title={\ding{45}\quad Exercice~5.9. Le dossier \niveauexpert}]
Établissez la liste des pièces que LiZbiZ doit tenir à jour pour prouver,
à tout moment, qu'elle respecte les articles~26, 29 et~30.
\end{boxexo}

\begin{boxexo}[title={\ding{45}\quad Exercice~5.10. L'assurance \niveauexpert}]
Rédigez le cahier des charges d'une assurance cyber pour LiZbiZ~: garanties,
exclusions à refuser, obligations de l'assuré, plafonds.
\end{boxexo}

\bigskip

\begin{boxinfo}
\textbf{Pour aller plus loin.}
\begin{itemize}
    \item Loi n°~2010/012 relative à la cybersécurité et la cybercriminalité
          au Cameroun~\cite{loi2010012}.
    \item Loi n°~2024/017 relative à la protection des données à caractère
          personnel~\cite{loi2024017}.
\end{itemize}
\end{boxinfo}

\begin{boxinfo}
\textbf{Transition.} Les fondements sont posés. La Partie~II apprend à
analyser les risques, en commençant par les identifier
(Chapitre~\ref{chap:identifier}).
\end{boxinfo}
